\begin{abstract}
\textbf{Background:} Dermatological deep learning models have been shown to exhibit performance disparities across skin tones, with diagnostic accuracy often degrading substantially for darker skin presentations \cite{Daneshjou2022}. Such biases threaten the equitable deployment of clinical decision support tools and may exacerbate existing health disparities in dermatology.

\textbf{Methods:} We evaluate melanoma classification bias on the Dermatology Data Initiative (DDI) dataset \cite{DDI2021}, which provides skin-tone annotations \cite{Fitzpatrick1988} enabling subgroup analysis. An EfficientNet-B0 model is pretrained on HAM10000 and fine-tuned on DDI. We then apply targeted synthetic data augmentation using color-space perturbations and geometric transforms (here, saved photometric/geometric transforms via Albumentations [color jitter, noise, elastic, rotations], not generative models [GANs/Diffusion]), with augmentation factors of 0$\times$, 2$\times$, 5$\times$, and 10$\times$. Stages 2 and 3 share identical hyperparameters (lr\,$=$\,$5{\times}10^{-5}$, batch\,$=$\,32, pos\_weight\,$=$\,dynamic) for controlled comparison. Synthetic images are generated exclusively from training-split dark-skin melanoma images ($n\,{=}\,29$, 10 variants each, 290 total). All experiments use seed 42; bootstrap 95\% confidence intervals (1{,}000 iterations) are reported.

\textbf{Results:} The pretrained baseline achieved overall AUROC 0.662 and Dark subgroup AUROC 0.588. Fine-tuning produced overall AUROC 0.634 and Dark AUROC 0.547, while adding 290 train-only synthetic dark melanomas produced overall AUROC 0.640 and Dark AUROC 0.481. The paired Dark AUROC delta for augmentation was $-0.066$ (95\% bootstrap interval $[-0.151,0.008]$, $p\,{=}\,0.962$; one-tailed directional, $H_1$: Dark improvement $>0$, $p$\,$=$\,proportion of bootstrap $\Delta \leq 0$). The ablation was non-monotonic, with the highest Dark AUROC at 2$\times$ (0.575) and the lowest at 10$\times$ (0.481).

\textbf{Conclusions:} In this single-seed study, simple photometric augmentation did not close the skin-tone performance gap and reduced Dark AUROC relative to fine-tuning. The findings are preliminary because the Dark test subgroup contains only 10 melanoma cases and there is no external validation. Larger multi-seed studies, independent cohorts, and clinically validated synthetic-image review are needed.
\end{abstract}
